\subsection{최종 답변의 자동 품질과 자원 사용 특성 (RQ3)}

\label{sec:ch4_4}

Jev의 자동 답변 지표는 Token F1이 1.12점(CI [+0.02, +2.27]점), LF-Eval이 2.58점(CI [+0.37, +4.79]점) 높았으나 Token F1의 개선은 경계적이었고, 입력량 절감은 관찰되지 않았다(조건별 값은 표~\ref{tab:ch4_4}). LF-Eval은 gold 조문의 사용을 채점 기준에 포함하므로 근거 회수 개선과 독립적인 증거가 아니다(지표 정의는 S1).

답변 생성 단계의 선택 문맥은 평균 203.67토큰에서 215.60토큰으로, 전체 API 입력은 2302.67토큰에서 2314.58토큰으로 늘었다. 문맥 길이의 문항별 차이 중앙값은 0이었고 선택 목록은 양쪽 모두 478개이므로, 이 증가는 선택 수가 아닌 선택 내용의 변화에 따른 것이다. LLM 선택 호출의 입력도 줄지 않았다(S6). 사용량은 신규 청구액이 아닌 조건별 요구량이다(S3).

선택 근거 수를 줄일 때의 자동 품질 변화를 탐색하기 위해, 저장된 최종 선택 근거에 scorer 점수 문턱값을 적용하는 선택 후 필터링을 수행하였다. 문턱값은 151개 의존 그룹의 5-fold 교차검증에서 학습 fold의 근거 F1 하락이 1점 이내인 후보 중 근거를 가장 적게 남기는 값으로 정하였고(결과 관찰 후 바꾼 격자 이력과 상세 절차는 S7), 유지 근거가 바뀐 문항(CE 10개, Jev 45개)만 같은 프롬프트와 모델로 답변을 재생성하여 평가하였다.

\begin{table}[H]
\centering
\caption{선택 후 필터링의 검증 fold 통합 결과와 답변 재평가(226문항). P·R은 근거 Precision·Recall, 완전 회수는 \%이다.}
\label{tab:ch4_4}
\scriptsize
\setlength{\tabcolsep}{1.5pt}
\begin{tabular}{@{}lrrrrrrr@{}}
\hline
조건 & 유지 근거 & F1 & P & R & 완전 회수 & Token F1 & LF-Eval \\
\hline
CE & 478 & 58.57 & 63.46 & 57.74 & 34.07 & 46.20 & 67.57 \\
CE+필터링 & 467 & 58.54 & 64.56 & 57.15 & 33.19 & 45.65 & 67.13 \\
Jev & 478 & 61.75 & 66.85 & 60.69 & 34.96 & 47.31 & 70.15 \\
Jev+필터링 & 426 & 61.43 & 69.80 & 58.26 & 33.19 & 47.03 & 69.13 \\
\hline
\end{tabular}
\end{table}

Jev+필터링(근거 89.12\% 유지)의 근거 F1 차이는 $-$0.32점(CI [$-$1.76, +1.07]점)으로 0 근처였으나, F1이 문턱값 선택 기준이었으므로 부분적으로 설계상의 귀결이다. Precision은 2.95점(CI [+1.05, +5.01]점) 높아졌고, Recall은 2.43점(CI [$-$3.81, $-$1.16]점), 완전 회수율은 1.77\%p(CI [$-$3.51, $-$0.43]\%p) 낮아졌다. 답변 지표의 차이는 Token F1 $-$0.28점(CI [$-$1.10, +0.44]점), LF-Eval $-$1.02점(CI [$-$2.23, +0.15]점)이었다. CE는 근거의 97.70\%를 유지하여 필터링 강도가 달랐으므로 조건 간 비교는 S7에 제시하였다. 사전에 정한 비열등성 마진이 없으므로 이 결과를 품질의 유지나 동등성으로 해석하지 않으며, 필터링에 따른 비용·지연 변화는 측정하지 않았다.
